% Assumptions
\newcommand{\RFassump}{
% assumptions for RCIF Adaptive Bounds!
\newcommand{\ASSUMPrcifconsisBoundedTau}{
\begin{assumption}
\label{assump:CuiAssump1}
There exists a fixed maximum follow-up time $0 < \tau < \infty$ and a constant $\tauUB \in (0, 1)$, such that $Pr(X_\tau \geq \tau | Z = z) \geq \tauUB$, for all $z \in \mathcal{Z}$.
\end{assumption}
}

\newcommand{\ASSUMPrcifconsisBoundedCovDistrO}{
\begin{assumption}
    \label{assump:CuiAssump2}
    The covariates $Z \in [0; 1]^d$ are distributed according to a density function $p(\cdot)$ satisfying $1/\assumpcovbound \leq p(z) \leq \assumpcovbound$, for all $z$ and some $\assumpcovbound \geq 1$.
\end{assumption}
}

\newcommand{\ASSUMPrcifconsisBoundedBelowKO}{
\begin{assumption}
\label{assump:CuiAssump3}
Assume that the minimum terminal node size $k$ is bounded below such that
\[
\lim_{n \to \infty} \frac{\log(n) \max\{\log(d), \log \log(n)\}}{k} = 0.
\]
\end{assumption}
}

\newcommand{\ASSUMPrcifconsisLipschitzCtsO}{
\begin{assumption}
\label{assump:CuiAssump4}
For any fixed time point $t$, the distribution function for failure from any cause $H(t | z)$ is $L_1$-Lipschitz continuous in terms of $z$, and the cause-$j$ specific hazard function $\lambda_j(t | z)$ is $L_2$-Lipschitz continuous in terms of $z$;
that is, 
\begin{equation*}
|H(t | z_1) - H(t | z_2)| \leq L_1 \norm{z_1 - z_2}
\text{ and }
|\lambda_j(t | z_1) - \lambda_j(t | z_2)| \leq L_2 \norm{z_1 - z_2}, 
\end{equation*}
where $|| \cdot ||$ is the Euclidean norm.
\end{assumption}
}

\newcommand{\ASSUMPrcifconsisHazardBoundO}{
\begin{assumption}
\label{assump:HazardBound}
Assume $\lambda_j(s|z)$ is bounded from $0$ to $\tau$ by some constant  $\assumphazbound$. Then, $\max_{s} \lambda_j(s|z) \leq \assumphazbound.$
\end{assumption}
}

\newcommand{\assumptionsRCIFadaptiveboundsO}{
%\ASSUMPrcifconsisBoundedTau
%\ASSUMPrcifconsisBoundedCovDistr
\ASSUMPrcifconsisBoundedBelowK
\ASSUMPrcifconsisLipschitzCts
\ASSUMPrcifconsisHazardBound
}

\newcommand{\assumpsForestconsistencyO}{
\begin{assumption}[Terminal node size, polynomial]
\label{assump:Cho1}
The minimum size \( n_{\text{min}} \) of the terminal nodes grows at the rate \( n_{\text{min}} \bowtie n^\beta \), \( \frac{1}{2} < \beta < 1 \), where \( a \bowtie b \) implies that both \( a = O(b) \) and \( b = O(a) \).
\end{assumption}

\begin{assumption}
[\( \alpha \)-regular and random-split trees]
\label{assump:Cho2}
Trees are \( \alpha \)-regular and random-split with a constant \( 0 < \psi < 1 \).
\end{assumption}

\begin{assumption}[Lipschitz continuous and bounded survival and censoring cumulative densities]
\label{assump:Cho3}
There exist constants \( \lipS \) and \( \lipG \) such that \( |\survfunc(t | h_1) - \survfunc(t | h_2)| \leq \lipS\|h_1 - h_2\| \) and \( |G(t | h_1) - G(t | h_2)| \leq \lipG \|h_1 - h_2\| \) for all \( h_1, h_2 \in \history \), \( t \in [0, \tau] \), where \( G \) is the censoring cumulative distribution function (CDF) and \( \survfunc(\tau^- \mid \histlower)G(\tau^- \mid \histlower) > c_1 \) for all \( \histlower \) and some constant \( c_1 > 0 \).
\end{assumption}

%% assumptions for RCIF Adaptive Bounds!
\newcommand{\ASSUMPrcifconsisBoundedTau}{
\begin{assumption}
\label{assump:CuiAssump1}
There exists a fixed maximum follow-up time $0 < \tau < \infty$ and a constant $\tauUB \in (0, 1)$, such that $Pr(X_\tau \geq \tau | Z = z) \geq \tauUB$, for all $z \in \mathcal{Z}$.
\end{assumption}
}

\newcommand{\ASSUMPrcifconsisBoundedCovDistrO}{
\begin{assumption}
    \label{assump:CuiAssump2}
    The covariates $Z \in [0; 1]^d$ are distributed according to a density function $p(\cdot)$ satisfying $1/\assumpcovbound \leq p(z) \leq \assumpcovbound$, for all $z$ and some $\assumpcovbound \geq 1$.
\end{assumption}
}

\newcommand{\ASSUMPrcifconsisBoundedBelowKO}{
\begin{assumption}
\label{assump:CuiAssump3}
Assume that the minimum terminal node size $k$ is bounded below such that
\[
\lim_{n \to \infty} \frac{\log(n) \max\{\log(d), \log \log(n)\}}{k} = 0.
\]
\end{assumption}
}

\newcommand{\ASSUMPrcifconsisLipschitzCtsO}{
\begin{assumption}
\label{assump:CuiAssump4}
For any fixed time point $t$, the distribution function for failure from any cause $H(t | z)$ is $L_1$-Lipschitz continuous in terms of $z$, and the cause-$j$ specific hazard function $\lambda_j(t | z)$ is $L_2$-Lipschitz continuous in terms of $z$;
that is, 
\begin{equation*}
|H(t | z_1) - H(t | z_2)| \leq L_1 \norm{z_1 - z_2}
\text{ and }
|\lambda_j(t | z_1) - \lambda_j(t | z_2)| \leq L_2 \norm{z_1 - z_2}, 
\end{equation*}
where $|| \cdot ||$ is the Euclidean norm.
\end{assumption}
}

\newcommand{\ASSUMPrcifconsisHazardBoundO}{
\begin{assumption}
\label{assump:HazardBound}
Assume $\lambda_j(s|z)$ is bounded from $0$ to $\tau$ by some constant  $\assumphazbound$. Then, $\max_{s} \lambda_j(s|z) \leq \assumphazbound.$
\end{assumption}
}

\newcommand{\assumptionsRCIFadaptiveboundsO}{
%\ASSUMPrcifconsisBoundedTau
%\ASSUMPrcifconsisBoundedCovDistr
\ASSUMPrcifconsisBoundedBelowK
\ASSUMPrcifconsisLipschitzCts
\ASSUMPrcifconsisHazardBound
}

\newcommand{\assumpsForestconsistencyO}{
\input{Chapters/02_Proj1/02_theorems/02_assumptions/Cho_assumptions1_5}
}

% assumptions for ITR consistency
\newcommand{\assumptionsITRconsistencyO}{
\begin{assumption}
\label{assump:boundeddensity}
Define the random function $B(a,Z) = \valueoneTRUE - (1-\tol)\valueonemaxTRUE$, $a \in \mathcal{A}$. Assume $B(a,Z)$ has continuous, bounded densities at 0, for each $a \in \mathcal{A}$.
\end{assumption}

\begin{assumption}
\label{assump:boundedV0}
    $\sup_{a,Z} \left\| \valueTRUE \right\|$ is bounded by some $\sevenbound < \infty$.
\end{assumption}
}

%\ASSUMPrcifconsisBoundedCovDistr
\begin{assumption}
    \label{assump:CuiAssump2}
    The covariates $Z \in [0; 1]^d$ are distributed according to a density function $p(\cdot)$ satisfying $1/\assumpcovbound \leq p(z) \leq \assumpcovbound$, for all $z$ and some $\assumpcovbound \geq 1$.
\end{assumption}

\begin{assumption}
\label{khbddtv}
The kernel \( \kernel \) for the random cumulative incidence forest has bounded total variation. Specifically, there exists a constant \( C \) such that the total variation of \( \kernel \) satisfies
\[
\int_{-\infty}^{\infty} \abs{\kernel^\prime(x)} \, dx \leq C,
\]
where \( \kernel^\prime(x) \) denotes the derivative of the kernel \( \kernel \).
\end{assumption}
}

% assumptions for ITR consistency
\newcommand{\assumptionsITRconsistencyO}{
\begin{assumption}
\label{assump:boundeddensity}
Define the random function $B(a,Z) = \valueoneTRUE - (1-\tol)\valueonemaxTRUE$, $a \in \mathcal{A}$. Assume $B(a,Z)$ has continuous, bounded densities at 0, for each $a \in \mathcal{A}$.
\end{assumption}

\begin{assumption}
\label{assump:boundedV0}
    $\sup_{a,Z} \left\| \valueTRUE \right\|$ is bounded by some $\sevenbound < \infty$.
\end{assumption}
}

% Assumptions for Forest Consistency
Below are assumptions used for consistency of both forests. Assumptions \ref{assump:Cho1}-\ref{assump:Cho3} are necessary and realistic for the tree-based estimator terminal node size and splitting rules. Assumption \ref{assump:Cho1} requires the terminal node size to be not only sufficiently large absolutely but also sufficiently small relative to the sample size. Assumption \ref{assump:Cho2} follows from Definition \ref{def:regsplotalphatree}. Assumption \ref{assump:Cho3} adds smoothness to the survival function and censoring cumulative distribution function. Assumption \ref{assump:CuiAssump2} follows from \cite{cui_tree_2017} and \cite{wager_adaptive_2016}, which allows for dependency between covariates. Thus, non-hyper-rectangular or categorical history spaces are allowed with this assumption. Assumption \ref{khbddtv} is necessary for the random cumulative incidence forest because there is non-differentiability of the treatment policy.
\assumpsForestconsistency
}

\newcommand{\ITRassump}{
% assumptions for RCIF Adaptive Bounds!
\newcommand{\ASSUMPrcifconsisBoundedTau}{
\begin{assumption}
\label{assump:CuiAssump1}
There exists a fixed maximum follow-up time $0 < \tau < \infty$ and a constant $\tauUB \in (0, 1)$, such that $Pr(X_\tau \geq \tau | Z = z) \geq \tauUB$, for all $z \in \mathcal{Z}$.
\end{assumption}
}

\newcommand{\ASSUMPrcifconsisBoundedCovDistrO}{
\begin{assumption}
    \label{assump:CuiAssump2}
    The covariates $Z \in [0; 1]^d$ are distributed according to a density function $p(\cdot)$ satisfying $1/\assumpcovbound \leq p(z) \leq \assumpcovbound$, for all $z$ and some $\assumpcovbound \geq 1$.
\end{assumption}
}

\newcommand{\ASSUMPrcifconsisBoundedBelowKO}{
\begin{assumption}
\label{assump:CuiAssump3}
Assume that the minimum terminal node size $k$ is bounded below such that
\[
\lim_{n \to \infty} \frac{\log(n) \max\{\log(d), \log \log(n)\}}{k} = 0.
\]
\end{assumption}
}

\newcommand{\ASSUMPrcifconsisLipschitzCtsO}{
\begin{assumption}
\label{assump:CuiAssump4}
For any fixed time point $t$, the distribution function for failure from any cause $H(t | z)$ is $L_1$-Lipschitz continuous in terms of $z$, and the cause-$j$ specific hazard function $\lambda_j(t | z)$ is $L_2$-Lipschitz continuous in terms of $z$;
that is, 
\begin{equation*}
|H(t | z_1) - H(t | z_2)| \leq L_1 \norm{z_1 - z_2}
\text{ and }
|\lambda_j(t | z_1) - \lambda_j(t | z_2)| \leq L_2 \norm{z_1 - z_2}, 
\end{equation*}
where $|| \cdot ||$ is the Euclidean norm.
\end{assumption}
}

\newcommand{\ASSUMPrcifconsisHazardBoundO}{
\begin{assumption}
\label{assump:HazardBound}
Assume $\lambda_j(s|z)$ is bounded from $0$ to $\tau$ by some constant  $\assumphazbound$. Then, $\max_{s} \lambda_j(s|z) \leq \assumphazbound.$
\end{assumption}
}

\newcommand{\assumptionsRCIFadaptiveboundsO}{
%\ASSUMPrcifconsisBoundedTau
%\ASSUMPrcifconsisBoundedCovDistr
\ASSUMPrcifconsisBoundedBelowK
\ASSUMPrcifconsisLipschitzCts
\ASSUMPrcifconsisHazardBound
}

\newcommand{\assumpsForestconsistencyO}{
\begin{assumption}[Terminal node size, polynomial]
\label{assump:Cho1}
The minimum size \( n_{\text{min}} \) of the terminal nodes grows at the rate \( n_{\text{min}} \bowtie n^\beta \), \( \frac{1}{2} < \beta < 1 \), where \( a \bowtie b \) implies that both \( a = O(b) \) and \( b = O(a) \).
\end{assumption}

\begin{assumption}
[\( \alpha \)-regular and random-split trees]
\label{assump:Cho2}
Trees are \( \alpha \)-regular and random-split with a constant \( 0 < \psi < 1 \).
\end{assumption}

\begin{assumption}[Lipschitz continuous and bounded survival and censoring cumulative densities]
\label{assump:Cho3}
There exist constants \( \lipS \) and \( \lipG \) such that \( |\survfunc(t | h_1) - \survfunc(t | h_2)| \leq \lipS\|h_1 - h_2\| \) and \( |G(t | h_1) - G(t | h_2)| \leq \lipG \|h_1 - h_2\| \) for all \( h_1, h_2 \in \history \), \( t \in [0, \tau] \), where \( G \) is the censoring cumulative distribution function (CDF) and \( \survfunc(\tau^- \mid \histlower)G(\tau^- \mid \histlower) > c_1 \) for all \( \histlower \) and some constant \( c_1 > 0 \).
\end{assumption}

%% assumptions for RCIF Adaptive Bounds!
\newcommand{\ASSUMPrcifconsisBoundedTau}{
\begin{assumption}
\label{assump:CuiAssump1}
There exists a fixed maximum follow-up time $0 < \tau < \infty$ and a constant $\tauUB \in (0, 1)$, such that $Pr(X_\tau \geq \tau | Z = z) \geq \tauUB$, for all $z \in \mathcal{Z}$.
\end{assumption}
}

\newcommand{\ASSUMPrcifconsisBoundedCovDistrO}{
\begin{assumption}
    \label{assump:CuiAssump2}
    The covariates $Z \in [0; 1]^d$ are distributed according to a density function $p(\cdot)$ satisfying $1/\assumpcovbound \leq p(z) \leq \assumpcovbound$, for all $z$ and some $\assumpcovbound \geq 1$.
\end{assumption}
}

\newcommand{\ASSUMPrcifconsisBoundedBelowKO}{
\begin{assumption}
\label{assump:CuiAssump3}
Assume that the minimum terminal node size $k$ is bounded below such that
\[
\lim_{n \to \infty} \frac{\log(n) \max\{\log(d), \log \log(n)\}}{k} = 0.
\]
\end{assumption}
}

\newcommand{\ASSUMPrcifconsisLipschitzCtsO}{
\begin{assumption}
\label{assump:CuiAssump4}
For any fixed time point $t$, the distribution function for failure from any cause $H(t | z)$ is $L_1$-Lipschitz continuous in terms of $z$, and the cause-$j$ specific hazard function $\lambda_j(t | z)$ is $L_2$-Lipschitz continuous in terms of $z$;
that is, 
\begin{equation*}
|H(t | z_1) - H(t | z_2)| \leq L_1 \norm{z_1 - z_2}
\text{ and }
|\lambda_j(t | z_1) - \lambda_j(t | z_2)| \leq L_2 \norm{z_1 - z_2}, 
\end{equation*}
where $|| \cdot ||$ is the Euclidean norm.
\end{assumption}
}

\newcommand{\ASSUMPrcifconsisHazardBoundO}{
\begin{assumption}
\label{assump:HazardBound}
Assume $\lambda_j(s|z)$ is bounded from $0$ to $\tau$ by some constant  $\assumphazbound$. Then, $\max_{s} \lambda_j(s|z) \leq \assumphazbound.$
\end{assumption}
}

\newcommand{\assumptionsRCIFadaptiveboundsO}{
%\ASSUMPrcifconsisBoundedTau
%\ASSUMPrcifconsisBoundedCovDistr
\ASSUMPrcifconsisBoundedBelowK
\ASSUMPrcifconsisLipschitzCts
\ASSUMPrcifconsisHazardBound
}

\newcommand{\assumpsForestconsistencyO}{
\input{Chapters/02_Proj1/02_theorems/02_assumptions/Cho_assumptions1_5}
}

% assumptions for ITR consistency
\newcommand{\assumptionsITRconsistencyO}{
\begin{assumption}
\label{assump:boundeddensity}
Define the random function $B(a,Z) = \valueoneTRUE - (1-\tol)\valueonemaxTRUE$, $a \in \mathcal{A}$. Assume $B(a,Z)$ has continuous, bounded densities at 0, for each $a \in \mathcal{A}$.
\end{assumption}

\begin{assumption}
\label{assump:boundedV0}
    $\sup_{a,Z} \left\| \valueTRUE \right\|$ is bounded by some $\sevenbound < \infty$.
\end{assumption}
}

%\ASSUMPrcifconsisBoundedCovDistr
\begin{assumption}
    \label{assump:CuiAssump2}
    The covariates $Z \in [0; 1]^d$ are distributed according to a density function $p(\cdot)$ satisfying $1/\assumpcovbound \leq p(z) \leq \assumpcovbound$, for all $z$ and some $\assumpcovbound \geq 1$.
\end{assumption}

\begin{assumption}
\label{khbddtv}
The kernel \( \kernel \) for the random cumulative incidence forest has bounded total variation. Specifically, there exists a constant \( C \) such that the total variation of \( \kernel \) satisfies
\[
\int_{-\infty}^{\infty} \abs{\kernel^\prime(x)} \, dx \leq C,
\]
where \( \kernel^\prime(x) \) denotes the derivative of the kernel \( \kernel \).
\end{assumption}
}

% assumptions for ITR consistency
\newcommand{\assumptionsITRconsistencyO}{
\begin{assumption}
\label{assump:boundeddensity}
Define the random function $B(a,Z) = \valueoneTRUE - (1-\tol)\valueonemaxTRUE$, $a \in \mathcal{A}$. Assume $B(a,Z)$ has continuous, bounded densities at 0, for each $a \in \mathcal{A}$.
\end{assumption}

\begin{assumption}
\label{assump:boundedV0}
    $\sup_{a,Z} \left\| \valueTRUE \right\|$ is bounded by some $\sevenbound < \infty$.
\end{assumption}
}

% Assumptions for ITR Consistency
Below are assumptions for the proposed estimator for consistency. We assume bounding in Assumptions \ref{assump:boundeddensity}-\ref{assump:boundedV0} to account for nonregularity of the value function. This is a necessary and standard assumption for non-regularity conditions.
\assumptionsITRconsistency
}